% Standalone problem document, generated from the adjacent README.md.
% Compile with XeLaTeX. Mathematical content is maintained in Markdown.
\documentclass[11pt,a4paper]{article}
\usepackage[fontsize=10.5pt]{scrextend}
\usepackage[margin=22mm,headheight=15pt,headsep=9mm,footskip=11mm]{geometry}
\usepackage{fontspec}
\setmainfont{texgyrepagella}[Extension=.otf,UprightFont=*-regular,BoldFont=*-bold,ItalicFont=*-italic,BoldItalicFont=*-bolditalic]
\setsansfont{texgyreheros}[Extension=.otf,UprightFont=*-regular,BoldFont=*-bold,ItalicFont=*-italic,BoldItalicFont=*-bolditalic]
\setmonofont{texgyrecursor}[Extension=.otf,UprightFont=*-regular,BoldFont=*-bold,ItalicFont=*-italic,BoldItalicFont=*-bolditalic,Scale=MatchLowercase]
\usepackage{amsmath,amssymb}
\usepackage{unicode-math}
\setmathfont{texgyrepagella-math.otf}
\usepackage{xcolor,microtype,enumitem,needspace,fancyhdr}
\usepackage{bookmark}
\definecolor{ink}{HTML}{17344B}
\definecolor{accent}{HTML}{197D80}
\definecolor{muted}{HTML}{596773}
\hypersetup{colorlinks=true,linkcolor=accent,urlcolor=accent,pdftitle={RA-18 - The Goreinov--Tyrtyshnikov--Zamarashkin conjecture on square-submatrix inverse norms},pdfauthor={Open Problems in Numerical Linear Algebra}}
\urlstyle{same}
\setlength{\parindent}{0pt}
\setlength{\parskip}{5pt plus 1pt minus 1pt}
\linespread{1.04}
\setlength{\emergencystretch}{3em}
\widowpenalty=10000
\clubpenalty=10000
\displaywidowpenalty=10000
\allowdisplaybreaks[1]
\setlist{nosep,leftmargin=1.5em}
\providecommand{\tightlist}{\setlength{\itemsep}{0pt}\setlength{\parskip}{0pt}}
\providecommand{\pandocbounded}[1]{#1}
\pagestyle{fancy}
\fancyhf{}
\fancyhead[L]{\sffamily\footnotesize\color{muted}OPEN PROBLEMS IN NUMERICAL LINEAR ALGEBRA}
\fancyhead[R]{\sffamily\bfseries\footnotesize\color{accent}RA-18}
\fancyfoot[L]{\parbox[b]{0.9\textwidth}{\sffamily\fontsize{8}{10}\selectfont\color{muted}Editorial ratings; a bounded search cannot certify that a problem remains unsolved.\\Literature check: 2026-09-11}}
\fancyfoot[R]{\sffamily\footnotesize\color{muted}\thepage}
\renewcommand{\headrulewidth}{0.3pt}
\renewcommand{\headrule}{\hbox to\headwidth{\color{accent}\leaders\hrule height \headrulewidth\hfill}}
\makeatletter
\renewcommand{\section}{\Needspace{5\baselineskip}\@startsection{section}{1}{0pt}{12pt plus 2pt minus 2pt}{4pt}{\sffamily\bfseries\large\color{ink}}}
\renewcommand{\subsection}{\Needspace{4\baselineskip}\@startsection{subsection}{2}{0pt}{9pt plus 1pt minus 1pt}{3pt}{\sffamily\bfseries\normalsize\color{ink}}}
\makeatother
\setcounter{secnumdepth}{0}
\begin{document}
{\sffamily\small\color{muted}Randomized and low-rank approximation\par}
\vspace{3pt}
{\raggedright\hyphenpenalty=10000\exhyphenpenalty=10000\sffamily\bfseries\fontsize{21}{24}\selectfont\color{ink}The Goreinov--Tyrtyshnikov--Zamarashkin conjecture on square-submatrix inverse norms\par}
\vspace{7pt}
{\sffamily\small\textbf{Difficulty:} challenging\quad\textbf{Importance:} interesting to the community\par}
{\sffamily\small\bfseries\color{accent}Status: Partially resolved\par}
\vspace{4pt}
{\color{accent}\hrule height 0.8pt}
\vspace{6pt}
\textbf{Rating rationale:} Improving the classical bound is challenging and would sharpen stability guarantees for skeleton/CUR approximation.

\hypertarget{problem-statement}{%
\subsection{Problem statement}\label{problem-statement}}

For integers $1\le r<n$, let $U\in\mathbb R^{n\times r}$ satisfy $U^TU=I_r$.
Write $U_I=U(I,:)$ for $I\subseteq\{1,\ldots,n\}$ with $|I|=r$, and let $\|\cdot\|_2$ denote the spectral norm.

\textbf{Conjecture} (Goreinov--Tyrtyshnikov--Zamarashkin {[}1, Eq. (2.6){]}).
Every such $U$ admits a nonsingular square submatrix $U_I$ satisfying
\[
\|U_I^{-1}\|_2\le\sqrt n.
\]

In the notation of {[}1{]}, this is $t(r,n)\le\sqrt n$, where
\[
t(r,n)=
\max_{\substack{U\in\mathbb R^{n\times r}\\U^TU=I_r}}
\qquad
\min_{\substack{I\subseteq\{1,\ldots,n\}\\|I|=r,\ \det(U_I)\ne0}}
\|U_I^{-1}\|_2.
\]
The minimum is over nonsingular submatrices; at least one exists because $\operatorname{rank}(U)=r$.

\hypertarget{known-bounds-and-cases}{%
\subsection{Known bounds and cases}\label{known-bounds-and-cases}}

The bound
\[
t(r,n)\le\sqrt{r(n-r)+1}
\]
is well known and follows from various arguments: maximal volume {[}1{]}, Bischof--Stewart QR (BSQR) {[}2, 3{]}, or volume sampling {[}4, 5{]} just to name a few.
Also, $t(r,n)\ge\sqrt n$ whenever $r+1$ divides $n$ {[}1, Lem. 2.2{]}, so the conjectured constant is sharp.

The real case $r=1$ is elementary.
Nesterenko {[}6{]} proved the real case $n=4$, $r=2$; Sengupta--Pautov {[}7{]} subsequently proved $r=2$ for all $n$.
Orthogonal completion gives $t(r,n)=t(n-r,n)$, leaving $3\le r\le n-3$ as the general open range.
These inverse norms control pseudoskeleton approximation errors {[}1, Thms. 3.1--3.2{]}.

\hypertarget{proposed-complex-extension}{%
\subsection{Proposed complex extension}\label{proposed-complex-extension}}

Define $t_{\mathbb C}(r,n)$ analogously using $U\in\mathbb C^{n\times r}$ with $U^*U=I_r$, where $*$ denotes conjugate transpose.
The minimization again runs only over nonsingular $U_I$.
The classical bound $\sqrt{r(n-r)+1}$ also holds over $\mathbb C$ {[}3{]}.

\textbf{Conjecture (proposed by the contributor).} There is an absolute constant $\alpha$, independent of $r$ and $n$, such that
\[
t_{\mathbb C}(r,n)\le\alpha\sqrt n,
\qquad 1\le r<n.
\]
The dependence on $r,n$ and the best possible $\alpha$ remain to be determined.
The constant $\alpha=1$ fails: $t_{\mathbb C}(2,4)=\sqrt{3+\sqrt3}>2$ {[}8, 9{]}.
For two columns, Nesterenko {[}9, Prop. 1{]} proves
\[
t_{\mathbb C}(2,n)\le c_2\sqrt n,
\qquad c_2=\left(2-\frac2{\sqrt3}\right)^{-1/2}\approx1.08766,
\]
with equality whenever $4\mid n$, so any universal $\alpha$ must satisfy $\alpha\ge c_2$.

\hypertarget{references}{%
\subsection{References}\label{references}}

\begin{enumerate}
\def\labelenumi{\arabic{enumi}.}
\tightlist
\item
  S. A. Goreinov, E. E. Tyrtyshnikov, and N. L. Zamarashkin, \emph{A theory of pseudoskeleton approximations}, Linear Algebra and its Applications \textbf{261} (1997), 1--21.
  \href{https://doi.org/10.1016/S0024-3795(96)00301-1}{Published paper}.
\item
  A. Damle, \emph{Computing Strong Rank-Revealing Factorizations for Matrices with Orthonormal Rows}, arXiv: 2607.13532v1 (2026).
  \href{https://arxiv.org/pdf/2607.13532v1}{Preprint}.
\item
  A. I. Osinsky, \emph{Close to optimal column approximation using a single SVD}, Linear Algebra and its Applications \textbf{725} (2025), 359--377.
  \href{https://doi.org/10.1016/j.laa.2025.07.016}{Published paper}.
\item
  H. Avron and C. Boutsidis, \emph{Faster subset selection for matrices and applications}, SIAM Journal on Matrix Analysis and Applications \textbf{34} (2013), 1464--1499.
  \href{https://doi.org/10.1137/120867287}{Published paper}.
\item
  A. Cortinovis and D. Kressner, \emph{Adaptive Randomized Pivoting for Column Subset Selection, DEIM, and Low-Rank Approximation}, SIAM Journal on Matrix Analysis and Applications \textbf{47} (2026), 25--47.
  \href{https://doi.org/10.1137/24M1719189}{Published paper}.
\item
  Y. Nesterenko, \emph{Submatrices with the best-bounded inverses: revisiting the hypothesis}, arXiv: 2303.07492 (2023; revised 2024).
  \href{https://arxiv.org/abs/2303.07492}{Preprint}.
\item
  R. Sengupta and M. Pautov, \emph{On the submatrices with the best-bounded inverses}, arXiv: 2604.05944v5 (2026).
  \href{https://arxiv.org/html/2604.05944v5}{Preprint}.
\item
  Y. Nesterenko, \emph{Submatrices with the best-bounded inverses: Studying $\mathbb R^{n\times2}$ and $\mathbb C^{n\times2}$}, arXiv: 2408.16631v1 (2024).
  \href{https://arxiv.org/html/2408.16631v1}{Preprint}.
\item
  Y. Nesterenko, \emph{Submatrices with the best-bounded inverses: an asymptotically tight upper bound for $\mathbb C^{n\times2}$}, arXiv: 2604.24087v1 (2026).
  \href{https://arxiv.org/html/2604.24087v1}{Preprint}.
\end{enumerate}

\hypertarget{status-check-2026-09-11}{%
\subsection{Status check --- 2026-09-11}\label{status-check-2026-09-11}}

Checked {[}1{]}--{[}9{]} and searched for later resolutions; no general real-case resolution or resolution of the proposed complex extension was found.
The real two-column result is a preprint and supports the partially resolved status.
\end{document}
